\section{설계 식의 적용 사례}
\label{sec:cases}
\subsection{EEG 창 처리와 모델 입력}
8개 전압 채널을 $250\,\mathrm{Hz}$로 획득하는 프로필을 선택한다.
Cyton의 기본 무선 경로가 이 채널ㆍ표본률 조합을 제공한다\cite{cyton}.
분석 그래프는 전압 단위 변환, 길이 63의 대칭 FIR, 길이 250ㆍhop 25의 창,
정규화 공분산, 접공간 특징, 고정 모델 추정으로 구성한다.
FIR 계수와 학습된 접공간 기준은 그래프 매개변수 식별자로 고정한다.

창의 갱신 간격과 첫 표본부터 마지막 표본까지의 시간은
\begin{equation}
 T_{\mathrm{hop}}=\frac{25}{250}=0.1\,\mathrm{s},\qquad
 T_{\mathrm{span}}=\frac{249}{250}=0.996\,\mathrm{s}
\end{equation}
이다. FIR의 신호 군지연은 $31/250=0.124\,\mathrm{s}$이다.
초기 유효성 정책은 62개 이전 표본을 채운 후 필터 출력을 유효로 판정한다.
창의 시작 위상을 $n_0=62$로 선택하면,
완전히 유효한 첫 창의 끝은 원천 index $62+249=311$에 위치한다.
첫 원천 표본의 시각을 0으로 두면 해당 끝점은 $1.244\,\mathrm{s}$이다.

이 그래프의 형식 연결은 표~\ref{tab:eegtypes}와 같다.
\begin{table}[tb]
\centering\small
\caption{EEG 사례의 형식과 갱신 규칙}
\begin{tabularx}{\linewidth}{@{}P{.20\linewidth}YY@{}}
\toprule 연산 결과 & 값의 형식 & 시간ㆍ상태 계약 \\\midrule
전압 표본 & $\mathbb R^{8}$, V & $250\,\mathrm{Hz}$, 시계 $\kappa$ \\
FIR 출력 & $\mathbb R^{8}$, V & 채널별 이전 62표본 상태 \\
창 & $\mathbb R^{8\times250}$, V & hop 25, 끝점 anchor \\
공분산 & $\mathrm{SPD}(8)$, $\mathrm V^2$ & 창당 1개, 정규화 식~\eqref{eq:spd} \\
접공간 특징 & $\mathbb R^{36}$, 무차원 & 기준 $R$과 모델 입력 스키마 \\
예측 & 선언한 class score 또는 확률 & 활성 모델 버전과 가용 시각 \\
\bottomrule
\end{tabularx}
\label{tab:eegtypes}
\end{table}

수신 블록의 마지막 표본을 기준으로 입구 지연 4 ms, 큐 대기 2 ms,
연산 3 ms, 출력 전달 1 ms를 실행 계약으로 선택하면 끝점 이후의 반응시간은
\begin{equation}
 D_{\mathrm{end}}=4+2+3+1=10\,\mathrm{ms}
\end{equation}
로 계산한다. 계약 레코드는 이 네 매개변수의 적용 환경과 증거 식별자를 가진다.
창 길이, 신호 군지연과 실행 반응시간은 각자의 필드로 유지한다.

\subsection{시계 보정 주기와 시간 오차}
시계 보정 직후의 오차 반경을 $\epsilon_0=0.2\,\mathrm{ms}$,
보정 간격 동안의 상대 속도 오차 상한을 $\rho$로 둔다.
잔차를 이 오차 반경에 포함하는 분석 모형은
\begin{equation}
 \epsilon(t)=\epsilon_0+\rho t
 \label{eq:drift_example}
\end{equation}
이다. $\rho=50\,\mathrm{ppm}$이면 $t=10\,\mathrm{s}$에서 $\epsilon=0.7\,\mathrm{ms}$,
$\rho=100\,\mathrm{ppm}$이면 같은 시각에서 $\epsilon=1.2\,\mathrm{ms}$를 얻는다.
양의 속도 오차 상한 $\rho>0$과 허용 반경 $\epsilon_\star$를 선택하면 보정 간격은
\begin{equation}
 T_{\mathrm{sync}}\le\frac{\epsilon_\star-\epsilon_0}{\rho},
 \qquad \epsilon_\star\ge\epsilon_0
\end{equation}
를 따른다. 그림~\ref{fig:clock}은 식~\eqref{eq:drift_example}의 계산값을 나타낸다.

\begin{figure}[tb]
\centering
\begin{tikzpicture}
\begin{axis}[width=.83\linewidth,height=58mm,
 xlabel={보정 이후 경과시간 (s)},ylabel={오차 반경 (ms)},
 xmin=0,xmax=20,ymin=0,ymax=2.4,xtick={0,5,10,15,20},ytick={0,.5,1,1.5,2},
 axis lines=left,tick align=outside,
 label style={font=\small},tick label style={font=\small},
 legend style={font=\small,draw=none,at={(.03,.97)},anchor=north west},
 legend cell align=left]
\addplot[black,solid,line width=.65pt,domain=0:20,samples=2]{0.2+0.05*x};
\addlegendentry{$\rho=50\,\mathrm{ppm}$}
\addplot[black,dashed,line width=.65pt,domain=0:20,samples=2]{0.2+0.1*x};
\addlegendentry{$\rho=100\,\mathrm{ppm}$}
\end{axis}
\end{tikzpicture}
\caption{시계 보정 간격과 오차 반경의 해석적 관계. 초기 반경 0.2 ms와 두 속도 오차 상한을 식~\eqref{eq:drift_example}에 대입하였다.}
\label{fig:clock}
\end{figure}

\subsection{고밀도 전기생리의 payload와 버퍼}
Neuropixels 1.0의 AP 스트림은 384채널을 $30\,\mathrm{kHz}$로 제공하고,
SpikeGLX의 아날로그 저장값은 signed 16-bit word를 사용한다\cite{neuropixels,spikeglx_metadata}.
이 AP 표본 배열의 payload 유량은
\begin{equation}
 r_{16}=384\times30\,000\times2=23.04\,\mathrm{MB/s}
\end{equation}
이다. 동일 표본 배열의 \code{f64} 표현은
\begin{equation}
 r_{64}=384\times30\,000\times8=92.16\,\mathrm{MB/s}
\end{equation}
를 사용한다. 여기서 MB는 $10^6$ byte이다. 각 값은 AP 표본 배열의 명시한 표현을 계량한다.
LFP, SY, 헤더와 재전송은 해당 경로의 추가 항으로 계량한다.

계산 모형은 기준 시계에서 균일한 $30\,000\,\mathrm{s}^{-1}$의 표본 발생을 사용한다.
한 블록을 채널당 30표본으로 구성하면 블록률은 $1000\,\mathrm{s}^{-1}$이고,
16-bit 블록의 크기는 $b=23\,040$ byte이다. 해당 FIFO 경계에
$R=50\,\mathrm{MB/s}$, $T=2\,\mathrm{ms}$의 서비스 곡선을 부여하면
\begin{equation}
 \begin{aligned}
 B_{\max}&=b+r_{16}T=69\,120\,\mathrm{byte}=67.5\,\mathrm{KiB},\\
 W_{\max}&=T+\frac bR=2.4608\,\mathrm{ms}.
 \end{aligned}
\end{equation}
블록 첫 표본부터 완성까지의 시간은 $29/30\,000\,\mathrm{s}$이다.
여기에 FIFO 지연을 합하면 첫 표본에서 선택한 출력 경계까지 $3.42747\,\mathrm{ms}$의
계산 상한을 얻는다. 서비스 계약은 블록 완성 경계를 기준으로 $R,T$를 지정한다.

연속 유입을 $T_{\mathrm{pause}}$ 동안 보관하는 값 배열의 용량은
$M=rT_{\mathrm{pause}}$로 계산한다. 그림~\ref{fig:memory}는 두 표현의 관계를 나타낸다.
\begin{figure}[tb]
\centering
\begin{tikzpicture}
\begin{axis}[width=.83\linewidth,height=58mm,
 xlabel={보관 시간 (ms)},ylabel={표본 배열 용량 (MB)},
 xmin=0,xmax=200,ymin=0,ymax=20,xtick={0,50,100,150,200},ytick={0,5,10,15,20},
 axis lines=left,tick align=outside,
 label style={font=\small},tick label style={font=\small},
 legend style={font=\small,draw=none,at={(.03,.97)},anchor=north west},
 legend cell align=left]
\addplot[black,solid,line width=.65pt,domain=0:200,samples=2]{23.04*x/1000};
\addlegendentry{16-bit: $23.04\,\mathrm{MB/s}$}
\addplot[black,dashed,line width=.65pt,domain=0:200,samples=2]{92.16*x/1000};
\addlegendentry{64-bit: $92.16\,\mathrm{MB/s}$}
\end{axis}
\end{tikzpicture}
\caption{384채널, 30 kHz AP 표본 배열의 보관 용량. 식 $M=rT_{\mathrm{pause}}$에 16-bit와 64-bit 표현의 유량을 대입하였다.}
\label{fig:memory}
\end{figure}

\subsection{EEG와 fNIRS의 시간 결합}
EEG 스트림은 전극 전압 스키마를, fNIRS 스트림은
$(s,d,\lambda,q,u)$의 광학 스키마를 유지한다. Aurora의 원시 광강도와 혈색소 변화량
출력은 서로 다른 명목 측정량을 갖는다\cite{nirx_formats}.
결합 연산자는 공통 시간 창 $[u_k,v_k)$와 각 입력의 시계 변환을 지정한다.
각 스트림에서 해당 창의 특징을 계산하고, 공통 frontier가 $v_k$에 도달하면
두 특징을 $(\varphi_{\mathrm{EEG}},\varphi_{\mathrm{fNIRS}})$의 곱 형식으로 출력한다.
후속 모델은 특징의 시간 지지, 표본 수, 유효성 마스크와 변환 계보를 함께 받는다.
